\documentclass[../../main.tex]{subfiles}% 注意这里的文件路径不能用 ./main.tex ,否则用latexmk编译子文件会报错
\graphicspath{{\subfix{./image/}}{\subfix{../../image/}}} % 指定图片引用路径,后续可以直接使用图片文件名
% 如果图片存放在上级目录,则使用 ../../image/ 作为备用路径

% 例如:
% \begin{figure}[H]
% \centering
% \includegraphics[width=0.3\textwidth]{图.png}
% \caption{}
% \label{figure:图}
% \end{figure}
% 注意:上述\label{}一定要放在\caption{}之后,否则引用图片序号会只会显示??.

\begin{document}

\section{不定积分的微分}

\begin{lemma}\label{lemma:引理变限积分的平均连续性}
设 $f\in L([a,b])$, 令
\begin{align*}
F_h(x) = \frac{1}{h}\int_x^{x+h} f(t)\mathrm{d}t
\end{align*}
(当 $x\notin [a,b]$ 时, 令 $f(x)=0$), 则有
\begin{align*}
\lim_{h\to 0}\int_a^b |F_h(x)-f(x)|\mathrm{d}x = 0.
\end{align*}
\end{lemma}
\begin{proof}
不妨设 $h>0$. 由于
\begin{align*}
F_h(x)-f(x) = \frac{1}{h}\int_x^{x+h} f(t)\mathrm{d}t - f(x) 
= \frac{1}{h}\int_0^h f(x+t)\mathrm{d}t - f(x) 
= \frac{1}{h}\int_0^h (f(x+t)-f(x))\mathrm{d}t,
\end{align*}
故知
\begin{align*}
\int_a^b |F_h(x)-f(x)|\mathrm{d}x \leqslant \int_{-\infty}^{+\infty} \left(\frac{1}{h}\int_0^h |f(x+t)-f(x)|\mathrm{d}t\right)\mathrm{d}x 
= \int_0^h \frac{1}{h}\mathrm{d}t \int_{-\infty}^{+\infty} |f(x+t)-f(x)|\mathrm{d}x.
\end{align*}
因为 $f\in L(-\infty,+\infty)$, 由\hyperref[theorem:积分的平均连续性]{积分的平均连续性}知, 对任给的 $\varepsilon>0$, 存在 $\delta>0$, 使得当 $|t|<\delta$ 时, 有
\begin{align*}
\int_{-\infty}^{+\infty} |f(x+t)-f(x)|\mathrm{d}x < \varepsilon,
\end{align*}
所以当 $h<\delta$ 时, 有
\begin{align*}
\int_0^h \frac{1}{h}\mathrm{d}t \int_{-\infty}^{+\infty} |f(x+t)-f(x)|\mathrm{d}x < \frac{\varepsilon}{h}\int_0^h \mathrm{d}x = \varepsilon.
\end{align*}
这说明, 对任给的 $\varepsilon>0$, 当 $|h|<\delta$ 时, 我们得到
\begin{align*}
\int_a^b |F_h(x)-f(x)|\mathrm{d}x < \varepsilon.
\end{align*}
\end{proof}

\begin{theorem}
设 $f\in L([a,b])$, 令
\begin{align*}
F(x) = \int_a^x f(t)\mathrm{d}t, \quad x\in[a,b],
\end{align*}
则
\begin{align*}
F'(x) = f(x), \quad \text{a.e. } x\in[a,b].
\end{align*}
即对于 $f\in L([a,b])$, 有
\begin{align*}
\frac{\mathrm{d}}{\mathrm{d}x}\int_a^x f(t)\mathrm{d}t = f(x), \quad \text{a.e. } x\in[a,b].
\end{align*}
\end{theorem}
\begin{proof}
令
\begin{align*}
F_h(x) = \frac{1}{h}\int_x^{x+h} f(t)\mathrm{d}t
\end{align*}
(当 $x\notin [a,b]$ 时, 令 $f(x)=0$). 因为 $F(x)$ 是几乎处处可微的, 所以不妨假定
\begin{align*}
\lim_{h\to 0} F_h(x) = g(x), \quad \text{a.e. } x\in[a,b],
\end{align*}
$g(x)$ 是 $[a,b]$ 上的可测函数. \cref{lemma:引理变限积分的平均连续性}指出, $F_h(x)$ 是平均收敛于 $f(x)$ 的 ($h\to 0$), 从而再由\hyperref[theorem:Fatou引理]{Fatou引理}得
\begin{align*}
\int_a^b |f(x)-g(x)|\mathrm{d}x = \int_a^b \lim_{h\to 0} |f(x)-F_h(x)|\mathrm{d}x \leqslant \varliminf_{h\to 0} \int_a^b |f(x)-F_h(x)|\mathrm{d}x = 0.
\end{align*}
这说明 $g(x)=f(x)$, a.e. $x\in[a,b]$.
\end{proof}

\begin{corollary}\label{corollary:可积函数的Lebesgue点}
若 $f\in L([a,b])$, 则对 $[a,b]$ 中几乎处处的点 $x$, 都有
\begin{align*}
\lim_{h\rightarrow 0} \frac{1}{h}\int_0^h{\left| f\left( x+t \right) -f\left( x \right) \right|\mathrm{d}t}=0 \iff \lim_{h\rightarrow 0} \frac{1}{h}\int_0^h{f\left( x+t \right) \mathrm{d}t}=f\left( x \right) .
\end{align*}
我们称满足上式的点 $x$ 为 $f(x)$ 的$\mathbf{Lebesgue}$\textbf{点}.
\end{corollary}
\begin{proof}
对于任意取定的 $r\in\mathbb{Q}\cap[a,b]$, 由\cref{lemma:引理变限积分的平均连续性}知
\begin{align*}
\lim_{h\to 0} \frac{1}{h}\int_x^{x+h} |f(t)-r|\mathrm{d}t = |f(x)-r|, \quad \text{a.e. } x\in[a,b].
\end{align*}
记 $[a,b]$ 中不满足上式的 $x$ 的全体为 $Z_r$, 则有 $m(Z_r)=0$. 又令
\begin{align*}
Z = \left(\bigcup_{r\in\mathbb{Q}\cap[a,b]} Z_r\right) \cup \{x\in[a,b]: |f(x)|=+\infty\},
\end{align*}
易知 $m(Z)=0$.

现在设 $x$ 为 $[a,b]$ 中不属于 $Z$ 的任一点, 则对任给 $\varepsilon>0$, 必有 $r\in\mathbb{Q}\cap[a,b]$, 使得
\begin{align*}
|f(x)-r|<\frac{\varepsilon}{3}.
\end{align*}
因为 $x\overline{\in}Z$, 所以存在 $\delta>0$, 当 $0<|h|<\delta$ 时, 有
\begin{align*}
\left|\frac{1}{h}\int_x^{x+h} |f(t)-r|\mathrm{d}t - |f(x)-r|\right| < \frac{\varepsilon}{3}\Longrightarrow \frac{1}{h}\int_x^{x+h}{\left| f\left( t \right) -r \right|\mathrm{d}t}<\left| f\left( x \right) -r \right|+\frac{\varepsilon}{3}.
\end{align*}
综上所述, 当 $0<|h|<\delta$ 时，我们有
\begin{align*}
\frac{1}{h}\int_x^{x+h}{\left| f\left( t \right) -f\left( x \right) \right|\mathrm{d}t}&\leqslant \frac{1}{h}\int_x^{x+h}{\left| f\left( t \right) -r \right|\mathrm{d}t}+\frac{1}{h}\int_x^{x+h}{\left| f\left( x \right) -r \right|\mathrm{d}t}
\\
&<\left| f\left( x \right) -r \right|+\frac{\varepsilon}{3}+\left| f\left( x \right) -r \right|<\frac{\varepsilon}{3}+\frac{\varepsilon}{3}+\frac{\varepsilon}{3}=\varepsilon .
\end{align*}
即得所证.再由
\begin{align*}
0\leqslant \lim\limits_{h\rightarrow 0^+} \frac{1}{h}\left| \int_0^h{f\left( x+t \right) -f\left( x \right) \mathrm{d}t} \right|\leqslant \lim\limits_{h\rightarrow 0^+} \frac{1}{h}\int_0^h{\left| f\left( x+t \right) -f\left( x \right) \right|\mathrm{d}t}=0,
\end{align*}
可得
\begin{align*}
\lim_{h\rightarrow 0^+} \frac{1}{h}\int_0^h{f\left( x+t \right) \mathrm{d}t}=\frac{1}{h}\int_0^h{f\left( x \right) \mathrm{d}t}=f\left( x \right) .
\end{align*}
\end{proof}

\begin{proposition}
设 $f\in L(\mathbb{R})$. 若对区间 $[a,b]$, 有
\begin{align*}
\int_a^b |f(x+h)-f(x)|\mathrm{d}x = o(|h|) \quad (h\to 0),
\end{align*}
则
\begin{align*}
f(x) = C \text{ (常数)}, \quad \text{a.e. } x\in[a,b].
\end{align*}
\end{proposition}
\begin{proof}
设 $x_1,x_2$ ($x_1<x_2$) 是 $f(x)$ 在 $(a,b)$ 中的 Lebesgue 点, 则由题设可知
\begin{align*}
\lim_{h\to 0} \left|\frac{1}{h}\int_{x_1}^{x_2} (f(x+h)-f(x))\mathrm{d}x\right| = 0.
\end{align*}
再利用\cref{corollary:可积函数的Lebesgue点}可得
\begin{align*}
&\left|\frac{1}{h}\int_{x_1}^{x_2} (f(x+h)-f(x))\mathrm{d}x\right|
= \left|\frac{1}{h}\int_{x_1+h}^{x_2+h} f(x)\mathrm{d}x - \frac{1}{h}\int_{x_1}^{x_2} f(x)\mathrm{d}x\right| \\
&= \left|\frac{1}{h}\int_{x_2}^{x_2+h} f(x)\mathrm{d}x - \frac{1}{h}\int_{x_1}^{x_1+h} f(x)\mathrm{d}x\right| 
\to |f(x_2)-f(x_1)|, \quad h\to 0,
\end{align*}
所以 $f(x_1)=f(x_2)$. 由于 $[a,b]$ 中几乎处处都是 $f(x)$ 的 Lebesgue点, 故 $f(x)$ 在 $[a,b]$ 上几乎处处等于一个常数.
\end{proof}

\begin{proposition}
设 $f\in L([a,b])$, 且令 $F(x)=\int_a^x f(t)\mathrm{d}t$ ($x\in[a,b]$), 则 $F\in  BV([a,b])$, 且
\begin{align*}
\bigvee_a^b(F) \leqslant \int_a^b |f(x)|\mathrm{d}x.
\end{align*}
\end{proposition}
\begin{proof}
对分划 $\Delta: a=x_0<x_1<\cdots<x_n=b$, 易知
\begin{align*}
\sum_{i=1}^n |F(x_i)-F(x_{i-1})| = \sum_{i=1}^n \left|\int_{x_{i-1}}^{x_i} f(t)\mathrm{d}t\right| 
\leqslant \sum_{i=1}^n \int_{x_{i-1}}^{x_i} |f(t)|\mathrm{d}t = \int_a^b |f(t)|\mathrm{d}t.
\end{align*}
由此即得 $\bigvee_a^b(F)\leqslant\int_a^b |f(x)|\mathrm{d}x$.
\end{proof}










\end{document}